\section{Domain Schema Design}
\label{sec:schema}

A well-specified schema is the contract between the LLM extraction agent and every downstream consumer of the knowledge graph.
The schema defines 17 entity types and 30 relation types, organized into functional categories that mirror the vocabulary of translational biomedicine.
The full schema is rendered as a linkage diagram in Figure~\ref{fig:schema}.

\subsection{Entity Type Taxonomy}

Entity types are grouped into six categories (Table~\ref{tab:entity-types}).
The taxonomy deliberately stays coarse-grained: each type maps to one or more
public ontologies (Section~\ref{sec:ontology-grounding}), and fine-grained
distinctions (e.g., kinase vs.\ phosphatase) are captured via ontology
identifiers rather than additional types.

\begin{table}[t]
\centering
\caption{Entity type taxonomy. Seventeen types organized into six functional categories.}
\label{tab:entity-types}
\begin{tabular}{lll}
\toprule
\textbf{Category} & \textbf{Entity Types} & \textbf{Count} \\
\midrule
Drug \& Chemistry      & \texttt{COMPOUND}, \texttt{METABOLITE}                                                                          & 2 \\
Molecular Biology      & \texttt{GENE}, \texttt{PROTEIN}, \texttt{PROTEIN\_DOMAIN}, \texttt{SEQUENCE\_VARIANT}, \texttt{CELL\_OR\_TISSUE} & 5 \\
Disease \& Phenotype   & \texttt{DISEASE}, \texttt{PHENOTYPE}, \texttt{ADVERSE\_EVENT}                                                    & 3 \\
Clinical               & \texttt{CLINICAL\_TRIAL}, \texttt{BIOMARKER}, \texttt{REGULATORY\_ACTION}                                       & 3 \\
Pathways \& Mechanisms & \texttt{PATHWAY}, \texttt{MECHANISM\_OF\_ACTION}                                                                & 2 \\
Context                & \texttt{ORGANIZATION}, \texttt{PUBLICATION}                                                                     & 2 \\
\bottomrule
\end{tabular}
\end{table}

\subsection{Relation Type Taxonomy}

The 30 relation types are organized into eight functional groups (Table~\ref{tab:relation-types}).
Most relations are directed; for example, \texttt{TARGETS} connects a \texttt{COMPOUND} subject to a \texttt{PROTEIN} object.
A single catch-all relation, \texttt{ASSOCIATED\_WITH}, absorbs associations that the model can identify but cannot classify into a more specific type, preserving recall at the cost of precision.

\begin{table}[t]
\centering
\caption{Relation type taxonomy. Thirty directed relation types in eight functional groups.}
\label{tab:relation-types}
\begin{tabular}{llr}
\toprule
\textbf{Group} & \textbf{Relation Types} & \textbf{Count} \\
\midrule
Drug--Target & \texttt{TARGETS}, \texttt{INHIBITS}, \texttt{ACTIVATES}, \texttt{BINDS\_TO}, \texttt{HAS\_MECHANISM} & 5 \\
Drug--Disease & \texttt{INDICATED\_FOR}, \texttt{CONTRAINDICATED\_FOR} & 2 \\
Drug--Clinical & \texttt{EVALUATED\_IN}, \texttt{CAUSES} & 2 \\
Drug--Drug & \texttt{DERIVED\_FROM}, \texttt{COMBINED\_WITH}, \texttt{INTERACTS\_WITH} & 3 \\
Molecular Biology & \makecell[tl]{\texttt{ENCODES}, \texttt{PARTICIPATES\_IN}, \texttt{IMPLICATED\_IN}, \texttt{CONFERS\_RESISTANCE\_TO},\\
\texttt{EXPRESSED\_IN}, \texttt{LOCALIZED\_TO}, \texttt{HAS\_VARIANT}, \texttt{HAS\_DOMAIN},\\
\texttt{METABOLIZED\_BY}, \texttt{PHOSPHORYLATES}, \texttt{FORMS\_COMPLEX\_WITH},\\
\texttt{REGULATES\_EXPRESSION}} & 12 \\
Biomarker & \texttt{PREDICTS\_RESPONSE\_TO}, \texttt{DIAGNOSTIC\_FOR} & 2 \\
Organizational & \texttt{DEVELOPED\_BY}, \texttt{PUBLISHED\_IN}, \texttt{GRANTS\_APPROVAL\_FOR} & 3 \\
Fallback & \texttt{ASSOCIATED\_WITH} & 1 \\
\bottomrule
\end{tabular}
\end{table}

\subsection{Ontology Grounding}
\label{sec:ontology-grounding}

Every extracted entity is grounded to one or more public ontologies at extraction time.
Table~\ref{tab:ontology-mapping} lists the primary mappings.
Grounding serves two purposes: it disambiguates synonyms (e.g., ``imatinib'' vs.\ ``Gleevec'' both resolve to ChEMBL941) and enables downstream federation with external databases such as Open Targets and ClinicalTrials.gov.

\begin{table}[t]
\centering
\caption{Primary ontology mappings for each entity type.}
\label{tab:ontology-mapping}
\begin{tabular}{ll}
\toprule
\textbf{Entity Type} & \textbf{Ontologies / Databases} \\
\midrule
\texttt{COMPOUND}         & ChEBI~\cite{chebi2016}, DrugBank, ChEMBL \\
\texttt{GENE}             & HGNC, NCBI Gene \\
\texttt{PROTEIN}          & UniProtKB \\
\texttt{DISEASE}          & MeSH~\cite{mesh2000}, Disease Ontology \\
\texttt{PATHWAY}          & KEGG, Reactome \\
\texttt{ADVERSE\_EVENT}   & MedDRA~\cite{meddra1999} \\
\texttt{SEQUENCE\_VARIANT}& ClinVar, dbSNP \\
\bottomrule
\end{tabular}
\end{table}

\subsection{Naming Conventions}

Consistent surface-form normalization is enforced at extraction time to reduce post-processing deduplication.
Drugs are recorded as International Nonproprietary Names (INN)---``pembrolizumab'' rather than the trade name ``Keytruda.''
Gene symbols follow HGNC conventions (\texttt{EGFR}, \texttt{TP53}, \texttt{BRCA1}).
Sequence variants use HGVS protein-level notation (e.g., BRAF V600E, KRAS G12C).
Diseases are normalized to MeSH~\cite{mesh2000} preferred terms, and adverse events to MedDRA~\cite{meddra1999} preferred terms.
These conventions are encoded directly in the extraction prompt (Section~\ref{sec:extraction}), so the LLM produces normalized names without a separate entity-linking pass.
